% !TEX root = CoeneDylanBP.tex
%%=============================================================================
%% Evaluatie en Resultaten
%%=============================================================================

\chapter{Evaluatie en Resultaten}
\label{ch:evaluatie}

Dit hoofdstuk bespreekt de resultaten van de ontwikkelde pipeline per component. Omdat er geen toegang was tot het HDTS-meetsysteem als referentie, ontbreekt kwantitatieve grondwaarheid (\textit{ground truth}) voor HD en ToF. De evaluatie is bijgevolg kwalitatief van aard: op basis van visuele observaties en interne consistentiechecks worden de sterktes en zwaktes van elk onderdeel in kaart gebracht. Aan het einde van dit hoofdstuk worden de resultaten getoetst aan de criteria uit de onderzoeksdoelstelling (Sectie~\ref{sec:onderzoeksdoelstelling}).

%----------------------------------------------------------------------

\section{Hoekpuntdetectie}
\label{sec:evaluatie-hoekpunten}

Het YOLO11n-pose model voor hoekpuntdetectie behaalt een hoge mAP50-score op de validatieset, wat wijst op een sterke generalisatie over de uiteenlopende camerastandpunten en -afstanden in het wedstrijdmateriaal. In de meeste frames detecteert het model de vier trampolinehoeken nauwkeurig en stabiel.

De voornaamste zwakte manifesteert zich op het moment van landing: wanneer de springer het doek raakt, zakt het trampolineoppervlak plaatselijk in. Dit veroorzaakt een vervorming van de zichtbare contour van de trampoline, waardoor het model de hoekpunten soms laat wegschuiven naar incorrect gedetecteerde posities. Omdat de volledige homografietransformatie op deze vier punten steunt, werkt elke afwijking hier rechtstreeks door in de berekende landingspositie.

De vergelijking tussen automatische detectie en manuele annotatie bevestigt dit: bij manueel vastgelegde hoekpunten is de landingspositiebepaling beduidend consistenter. Manuele annotatie fungeert daarmee als bovengrens voor de haalbare nauwkeurigheid van het systeem in de huidige opzet.

%----------------------------------------------------------------------

\section{Homografietransformatie}
\label{sec:evaluatie-homografie}

De homografietransformatie converteert pixelco{\"o}rdinaten correct naar metrische posities voor punten die dicht bij de camera liggen, d.w.z. langs de links-rechts as van het camerastandpunt. Bij visuele verificatie (het projecteren van bekende markeringen op de mat terug naar beeldco{\"o}rdinaten) bedraagt de afwijking langs deze as slechts enkele centimeters.

De diepte-as (voor-achter vanuit het camerastandpunt) presteert beduidend minder nauwkeurig. Een homografie is algebraisch gevoelig voor kleine fouten in de hoekpunten, en dit effect is het sterkst langs de as die het verst van de camera verwijderd is. Een kleine fout in de gedetecteerde hoekpositie leidt langs de diepte-as tot een grotere metrische afwijking dan langs de parallelle as. Voor HD-zones die voornamelijk langs de links-rechts as worden onderscheiden, is de nauwkeurigheid aanvaardbaar; voor zones die langs de diepte-as worden bepaald, is de onzekerheid groter.

%----------------------------------------------------------------------

\section{Landingsmomentdetectie}
\label{sec:evaluatie-landingsmoment}

De binaire classifier onderscheidt de klassen \textit{Landed} en \textit{Not Landed} met hoge betrouwbaarheid. Visuele controle van gedetecteerde landingsframes toont aan dat het model de overgang van vlucht naar contact consistent en correct identificeert, zelfs bij varierende belichting en camerastandpunten.

Een systematisch effect dat optreedt is het \textit{eerste-frame probleem}: op het moment dat de classifier voor het eerst \textit{Landed} detecteert, heeft het trampolinedoek al licht ingezakt onder het gewicht van de springer. De springer staat op dit frame dus iets lager dan de werkelijke landingspositie, wat een kleine systematische fout introduceert in de positiebepaling. Dit effect is inherent aan de frame-gebaseerde aanpak en kan niet worden ge{\"e}limineerd zonder hogere framerate of sensorfusie.

%----------------------------------------------------------------------

\section{Landingspositiebepaling}
\label{sec:evaluatie-landingspositie}

De finale SAM-2-gebaseerde aanpak levert visueel plausibele segmentatiemaskers op voor de springer in het landingsframe. Het laagste pixel van het masker correspondeert in de meeste gevallen goed met de voetpositie van de springer op het doek.

Twee terugkerende problemen werden waargenomen. Ten eerste interfereren \textit{spotters} die dicht bij de trampoline staan: wanneer de YOLO-detector de begeleider als springer identificeert, genereert SAM-2 een masker van de verkeerde persoon, wat resulteert in een volledig foutieve landingspositie. Dit probleem treedt echter minder frequent op dan bij de pose estimation-aanpak, omdat de bounding box van YOLO explicieter de persoon op het trampolinedoek selecteert.

Ten tweede beïnvloedt de diepte-as beperking van de homografie (Sectie~\ref{sec:evaluatie-homografie}) ook de landingspositiebepaling: zelfs wanneer het contactpunt correct als pixel wordt geïdentificeerd, vertaalt een onnauwkeurige homografie dit naar een foutieve metrische positie langs de diepte-as.

%----------------------------------------------------------------------

\section{ToF-schatting}
\label{sec:evaluatie-tof}

De ToF-schatting via frame-telling levert intern consistente resultaten op: de berekende vluchttijden per sprong vertonen een realistisch patroon, met kortere intervallen voor lagere sprongen en langere intervallen voor hogere sprongen.

De inherente fps-begrenzing blijft echter de voornaamste beperking. Bij 25 fps bedraagt de onzekerheid per gedetecteerd landingsmoment 40\,ms. Over een volledige reeks van tien sprongen, met twintig landingsmomenten, loopt de cumulatieve onzekerheid op tot maximaal 800\,ms. De eis uit de onderzoeksdoelstelling --- een MAE kleiner dan 0,2 seconden --- kan bij 25 fps niet worden gegarandeerd zonder een betere landingsdetectie op subframe-niveau, wat niet mogelijk is zonder hogere framerate of een aanvullende sensor.

%----------------------------------------------------------------------

\section{Globale beoordeling}
\label{sec:evaluatie-globaal}

De ontwikkelde pipeline toont aan dat een low-cost computer vision-systeem de fundamentele componenten van HD- en ToF-meting kan realiseren. De totale hardwarekost van de gebruikte opstelling (een standaard camera en een laptop voor inferentie) blijft ruimschoots onder de gestelde grens van \euro{}500, waarmee het eerste criterium uit de onderzoeksdoelstelling wordt behaald.

Voor de classificatienauwkeurigheid van HD-zones kan geen kwantitatieve uitspraak worden gedaan bij gebrek aan grondwaarheid. Kwalitatief toont de pipeline correcte landingsposities aan bij manuele annotatie van de hoekpunten, en redelijke posities bij automatische detectie voor frames zonder sterke doekvervorming. De nauwkeurigheid langs de links-rechts as is aanvaardbaar; langs de diepte-as is verdere verbetering nodig.

De ToF-nauwkeurigheid is bij 25 fps niet voldoende om de gestelde MAE-eis te halen over een volledige reeks. Een camera met hogere framerate zou deze beperking grotendeels oplossen. De resultaten bevestigen dat het concept technisch haalbaar is, maar dat verdere iteraties noodzakelijk zijn om de vooropgestelde criteria volledig te behalen.
